\documentclass[10pt]{article}
\usepackage[a4paper,margin=25mm]{geometry}
\usepackage[T1]{fontenc}
\usepackage{lmodern,microtype}
\usepackage{amsmath,amssymb,amsthm,booktabs,longtable,hyperref}
\newtheorem{lemma}{Lemma}
\newcommand{\Seq}[2]{\langle #1,#2\rangle}
\newcommand{\CB}{\operatorname{break}}
\title{Source restructuring behind the 239-state ZFC candidate}
\author{Research supplement}
\date{19 October 2026}
\begin{document}
\maketitle

\section{What the source optimization contributes}
The final literal machine in \texttt{../candidate239/zfc\_boundary239.tm}
has 239 working states, a separate halt state, two symbols, and no stationary
moves. This note explains the source and encoding part of its construction.
The logical calibration to consistency and the complete quantified runtime
check are separate components of the accompanying report.

The inherited 240-state source has 6,455 instruction leaves and 225 distinct
binary subtrees. A source subtree is either an instruction or a pair
$\Seq{P}{Q}$ interpreted as sequential composition with explicit encoded
control exits. Increment and decrement instructions act on natural-valued
registers. A zero decrement follows a break edge; a positive decrement
follows the ordinary successor. Explicit break and continue instructions
count, respectively, zero and one ancestors in the binary instruction path.

A finite source control-flow graph can therefore be checked exactly. Give
each instruction its left-to-right ordinal, record its operation and argument,
and record the target ordinal for each possible exit. If two source trees
have the same ordered instructions and the same target ordinals, the identity
map on ordinals is a register-machine bisimulation. This statement holds for
all natural-valued register assignments, including computations that diverge.
It is stronger than agreement on finitely many test inputs.

\section{Associativity and the final break}
A subtree is syntactically closed when none of its local break or continue
instructions needs an ancestor outside that subtree, and it has no halt
instruction. For closed $P,Q,R$, the standard associativity transformation
\[
 \Seq{P}{\Seq{Q}{R}} \quad\longleftrightarrow\quad
 \Seq{\Seq{P}{Q}}{R}
\]
preserves both normal behavior and divergence. In the current source, the
initial useful reassociation is
\[
 \Seq{p}{\Seq{d}{h}}
 \quad\longmapsto\quad
 \Seq{\Seq{p}{d}}{h},
\]
where $p$ is the PUSH routine, $d$ copies register D, and $h$ is PUSH followed
by a copy of register H. The pair $\Seq{p}{d}$ already occurs elsewhere,
so the source DAG loses one node.

This reduction is insufficient at the tape level. Six decrement instruction
paths acquire a run of seven zeros, while the inherited backward scanner
requires the source's six-zero bound. Moreover, state 218 uses the old
reassociated decoder as a backward routing helper. The direct 239-state
attempt consequently fails 702 of 8,386 finite instruction guards, despite
having the same exact initial blank-tape prefix.

A beam of further closed reassociations removes five of the six bad paths
without increasing the 224-node source count. One bad path remains, at the
first PUSH decrement of the AX branch. The following extension removes it.

\begin{lemma}[Right association across a final break]
Let $P$ and $Q$ be syntactically closed and let $k\geq 1$. Then
\[
 \Seq{\Seq{P}{Q}}{\CB k}
 \quad\longleftrightarrow\quad
 \Seq{P}{\Seq{Q}{\CB k}}
\]
preserves the register computation and the external break target.
\end{lemma}
\begin{proof}
Each side executes $P$ followed by $Q$. If either diverges, both sides diverge
with the same register computation. Otherwise both reach the same registers
at the break. The relative path of that break changes from $1$ to $11$.
No zero ancestor is added or removed, so the counted external break target
is unchanged. Since $P$ and $Q$ are closed, their own instruction exits are
unaffected by changing their enclosing parentheses.
\end{proof}

The final AX subtree has the form
\[
 \Seq{\Seq{\mathrm{GEN}}{\Seq{\mathrm{AXrest}}{\mathrm{clear}}}}{\CB 1}.
\]
All three displayed computational routines are closed. Applying the lemma
twice, first at instruction prefix \texttt{11101001001} and then at the
resulting prefix \texttt{111010010011}, gives
\[
 \Seq{\mathrm{GEN}}{
   \Seq{\mathrm{AXrest}}{\Seq{\mathrm{clear}}{\CB 1}}}.
\]
The resulting source has 224 distinct subtrees, 6,455 leaves, maximum path
length 48, and maximum zero run 6. The exact source-graph check verifies every
ordered instruction and every conditional or explicit-control target. The
lemma explains why the former all-three-chunks-closed search rule missed the
useful transformation; the finite graph certificate checks the actual full
program independently.

\section{Changing the decoder-sharing choice}
A reduced source DAG is not yet a smaller correct Turing machine. The
inherited compiler shares some unequal source subtrees by using a bit already
present immediately left of the head. Reparenthesizing a source can change
the available tag or an intermediate route. This is the second obstruction
in the direct reduction.

Let $p$ denote PUSH, $f$ denote the source routine fn45, and $i$ denote
incrementing logical register 1. Their entry states are, respectively,
23, 105, and 70 in the inherited labels. The relevant source nodes are
\[
 F=\Seq{p}{i},\qquad G=\Seq{f}{i},\qquad H=\Seq{f}{f}.
\]
The old decoder represents $F$ and $G$ together at state 218, according to
whether the preceding tape bit is 0 or 1. State 179 separately represents $H$
and is entered after a 0 bit. The source reassociation changes a helper used
by the shared $F/G$ zero branch.

The successful repair represents $F$ explicitly at state 223:
\[
 223/0=0R23,\qquad 223/1=1R70,\qquad 237/0=0R223.
\]
All remaining entries into 218 now follow a transition that writes 1 and
moves right; all entries into 179 follow a transition that writes 0 and
moves right. These are global incoming-edge properties, so their invariants
hold from the common initial state on arbitrary tapes.

The construction then shares $G$ with $H$ instead of sharing $G$ with $F$.
Replace 179 by
\[
 179/0=0R105,\qquad 179/1=1L215,
\]
redirect entries into 218 to 179, and remove 218. The common first child is
$f$. The second child is selected by a route through existing states:
\begin{center}
\begin{tabular}{clc}
\toprule
Preceding bit & New read-1 path & Endpoint \\
\midrule
0 & $179\to215\to183\to213\to237\to105$ & $(105,h+1)$ \\
1 & $179\to215\to233\to70$ & $(70,h+1)$ \\
\bottomrule
\end{tabular}
\end{center}
Every step in these paths preserves the scanned tape bit. Only positions
$h-1$ and $h$ are read, and the complete tape is unchanged at the endpoint.
On read 0 the new direct transition to 105 matches old state 179, or replaces
old state 218's three-step route through $86/1$ and $99/0$.

There are four local cases: each of the two old states, under its incoming
tag, and each read symbol. The independent checker
\texttt{../candidate239/build\_candidate239.py} verifies all four by refusing
any exterior tape read. A positive finite-block simulation then gives
halting equivalence for the repaired 240-state reference and the final
239-state table. The companion fold note gives the complete connection to
the separately validated reference table.

Thus the useful operation is a change in the choice of shared decoder
nodes. The state removed by source reparenthesization is retained as a net
saving because a newly favorable pair of unequal source programs can share
one state using an inherited tape bit.

\section{Verification and provenance}
The canonical final table has 478 rows, start state 0, dispatcher 84, and
separate halt state 239. Its SHA-256 is
\begin{quote}\small\ttfamily
666f5bf1370b6b44ff4e614e4b194bc5be6b4c621914842ad7e344d6570a544e
\end{quote}
The inherited-label table has SHA-256
\begin{quote}\small\ttfamily
d061770f487937f9da652919c9bc9bb79bc9bf137ca41dd7a08eb5ecda8e7719.
\end{quote}
Canonicalization is a literal action-preserving bijection of all working and
halt labels. Sparse and dense interpreters, including a separately relabeled
dense run, agree on the exact 409,691-step blank prefix: the head is at $-5$,
the dispatcher is active, and exactly 351 tape cells contain one. The root
register values are $(3,0,0,0,3,0,0,0,0)$ with 167 spare unary values all equal
to one. No merged or split state is visited during this prefix.

The final source certificate's JSON bytes have SHA-256
\begin{quote}\small\ttfamily
7fd152d1f5371d2a561d4ad1870d241683b010ce7b7a27399b95363a4537cfb6.
\end{quote}
The quantified source-runtime checker reports all 8,386 guards verified with
2,247,513 symbolic rules. Those checks quantify over the register lengths
and the permitted runtime frame parameters; they are distinct from the
finite sanity screens used to reject preliminary candidates.

\section{Overlap and unsuccessful experiments}
The initial closed reassociation already appears as an \texttt{obvious244}
negative experiment in the inherited companion's independent runtime script.
It is not claimed as novel. The contribution of the present construction is
its completed combination with the boundary-aware source rotations and the
new sharing choice at 179/218, yielding a correct literal table below 240.
This is not a claim of exhaustive literature priority.

Further arithmetic loop rotations were also explored. If a transfer loop
adds a vector $V$ to registers other than its counter, and $A$ is a subvector
of $V$, then a loop followed by $A$ can be rotated into an initial $A$, a
decrement test, the remaining vector $V-A$, and a repeat. For counter value
$n$, both forms add $nV+A$. In this source the tested contextual versions
added four to seven distinct source nodes because the original shared loops
remained needed elsewhere. Thirty such variants did not give a smaller
literal candidate.

A separate decoder-aware beam tried to manufacture a zero-cost helper for
the original state-218 route while preserving another projection used by
states 243 and 220. It examined 930,448 generated alternatives over 20 levels
with width 240 and found no improvement under its bounded settings. An
unconstrained source result reached 224 nodes with the desired helper but
lost another inherited contextual sharing and failed 256 finite instruction
guards. These are useful cost-model counterexamples, not lower bounds.

The full source and literal certificates remain essential: a small source
DAG, a successful blank prefix, and finite runtime tests each establish only
part of the needed equivalence.
\end{document}
